\documentclass[11pt]{article}
\usepackage{amsmath,amsthm,amssymb}
\usepackage[margin=1in]{geometry}
\usepackage{hyperref}
\usepackage{enumitem}

\newtheorem{theorem}{Theorem}
\newtheorem{lemma}[theorem]{Lemma}
\newtheorem{corollary}[theorem]{Corollary}
\newtheorem{proposition}[theorem]{Proposition}
\theoremstyle{definition}
\newtheorem{definition}[theorem]{Definition}
\newtheorem{remark}[theorem]{Remark}

\newcommand{\hatl}{\widehat\lambda}
\newcommand{\SYT}{\mathrm{SYT}}

\title{An explicit column-descent-free SYT construction:\\
proof of Conjecture~4.6 of the converse-per-SYT-positivity paper}
\author{Clio}
\date{16 May 2026}

\begin{document}
\maketitle

\begin{abstract}
Let $\hatl = (\hatl_0, \ldots, \hatl_{L-1})$ be a partition with
$L \ge 2$ rows, each of length $\hatl_r \ge 2$.  Let $q \ge 1$ and
set $\lambda = (\hatl, 1^q)$.  Write $|\hatl|$ for the total number
of cells of $\hatl$, set $M = |\hatl| - 2L$ and $n = |\lambda| = 2L+M+q$.
Assume $\tau(\hatl) := \max(0, q + 2L - 1 - |\hatl|) = 0$, equivalently
$q \le M+1$.

We give an explicit deterministic construction $T^*$ of a standard
Young tableau of shape $\lambda$ such that
\begin{enumerate}[label=(\roman*)]
\item $T^*(0,0) = 1$ and $T^*(0,1) = 2$;
\item no $j \in \{1,\ldots,n-1\}$ has $T^*(j+1)$ directly below
$T^*(j)$ (the tableau is \emph{column-descent-free}).
\end{enumerate}
This proves Conjecture~4.6 of \cite{ConversePerSYT}, closing the
remaining structural gap in the converse direction of the
trace-vanishing equivalence.

The construction has a simple two-band form: place values $1,3,\ldots,2k-1$
in the top $k$ cells of column~$0$, where $k$ is the number of rows of
$\hatl$ of length at least~$3$; place the values $n, n-2, \ldots$ in the
remaining $L+q-k$ cells of column~$0$ from bottom to top, again spaced
by~$2$; then fill the rest of the diagram row-by-row in order with the
remaining values.  The verification consists of three elementary
inequalities, each tight.
\end{abstract}

\section{Setup and statement}\label{sec:setup}

Throughout we use $0$-indexed coordinates: row $0$ is the topmost row,
column $0$ the leftmost column.  A cell is written $(r,c)$.

Let $\hatl = (\hatl_0, \ldots, \hatl_{L-1})$ with $L \ge 2$ and
$\hatl_r \ge 2$ for all $r$; write $|\hatl| = \sum_r \hatl_r$, set
$M := |\hatl| - 2L$, and put $\lambda := (\hatl, 1^q)$ for some
$q \ge 1$.  Then $n := |\lambda| = 2L + M + q$, and the threshold
condition $\tau(\hatl) = 0$ reads $q \le M+1$.

We use $k$ throughout for the number of \emph{long} rows of $\hatl$:
\[
   k \;:=\; \#\{\, r \in \{0,\ldots,L-1\} : \hatl_r \ge 3 \,\}.
\]
Since $\hatl$ is weakly decreasing with each row $\ge 2$, the rows
$0, 1, \ldots, k-1$ have length $\ge 3$ and the rows $k, k+1, \ldots,
L-1$ have length exactly~$2$.

\begin{definition}[Column descent]\label{def:col-descent}
A standard Young tableau $T$ of shape $\lambda$ has a \emph{column
descent at $j$} if the cell of $j+1$ is directly below the cell of
$j$, i.e.\ $T^{-1}(j) = (r,c)$ and $T^{-1}(j+1) = (r+1, c)$ for some
$(r,c)$.  We say $T$ is \emph{column-descent-free} if no $j$ is a
column descent.
\end{definition}

\begin{theorem}[Main]\label{thm:main}
Under the hypotheses above, there exists a column-descent-free standard
Young tableau $T^*$ of shape $\lambda$ with $T^*(0,0) = 1$ and
$T^*(0,1) = 2$.
\end{theorem}

We prove the theorem by writing down an explicit $T^*$ and verifying
that it has the desired properties.

\section{The construction $T^*$}\label{sec:construction}

\paragraph{Column-$0$ values.}  Define for $r = 0, 1, \ldots, L+q-1$:
\begin{equation}\label{eq:vr}
  v_r \;:=\;
  \begin{cases}
     2r+1, & 0 \le r < k,\\
     n - 2(L+q-1-r), & k \le r \le L+q-1.
  \end{cases}
\end{equation}
Set $T^*(r, 0) := v_r$ for each $r \in \{0, 1, \ldots, L+q-1\}$.

\paragraph{Remaining cells.}  The cells of $\lambda$ not in column $0$
form a subdiagram $\hatl' := (\hatl_0 - 1, \hatl_1 - 1, \ldots,
\hatl_{L-1} - 1)$, with each row of length $\ge 1$.  Let
\[
   w_1 < w_2 < \cdots < w_{L+M}
\]
be the elements of $\{1, 2, \ldots, n\} \setminus \{v_0, \ldots,
v_{L+q-1}\}$ in increasing order.  Fill $\hatl'$ in row-superstandard
order: in row $r$ of $\hatl'$, place the values
$w_{a_r + 1}, w_{a_r + 2}, \ldots, w_{a_r + \hatl_r - 1}$ at columns
$1, 2, \ldots, \hatl_r - 1$ of $\lambda$, where
\[
   a_r \;:=\; \sum_{i < r} (\hatl_i - 1).
\]

The cells of $\hatl'$ are filled by $L+M$ values in row-major order;
together with the $L+q$ column-$0$ values $v_r$, all $n$ values are
placed.

\paragraph{Sectional structure of the $w_i$.}  We describe the
sequence $w_1 < w_2 < \cdots < w_{L+M}$ by three sections.

For $k \ge 1$: the values $\{1, 3, \ldots, 2k-1\}$ are in column $0$,
so the values $2, 4, \ldots, 2k-2$ are the smallest non-column-$0$
values; then there is an uninterrupted stretch of integers
$2k, 2k+1, \ldots, v_k - 1$ all non-column-$0$; finally, in the range
$[v_k, n]$, every other integer is in column $0$, leaving
$v_k + 1, v_k + 3, \ldots, n-1$ in $\hatl'$.  Explicitly:
\begin{equation}\label{eq:sections}
\begin{aligned}
\text{Section S1: }& w_i = 2i \text{ for } 1 \le i \le k-1\\
\text{Section S2: }& w_i = i + k \text{ for } k \le i \le k + (M - q + 1)\\
\text{Section S3: }& w_i = v_k + 2(i - (k + M - q + 1)) - 1\\
   & \quad\text{for } k + M - q + 2 \le i \le L + M.
\end{aligned}
\end{equation}
Section S1 has $k-1$ terms (gap $2$), Section S2 has $M-q+2$ terms
(gap $1$ --- the only consecutive integers), Section S3 has
$L+q-k-1$ terms (gap $2$).  Total: $(k-1) + (M-q+2) + (L+q-k-1) =
L + M$, as required.

For $k = 0$: every row of $\hatl$ has length $2$, so $|\hatl| = 2L$,
$M = 0$, and tau-zero forces $q = 1$; then $n = 2L+1$, the
column-$0$ values are $\{1, 3, \ldots, 2L+1\}$, and the non-column-$0$
values are $\{2, 4, \ldots, 2L\}$.  Sections S1 and S2 are empty, and
S3 has all $L = L+M$ terms, with $w_i = 2i$ for $1 \le i \le L$
(consistent with formula \eqref{eq:sections} for $i \in \{1, \ldots,
L\}$ in S3).  In particular, the sequence $w$ has no consecutive
integers when $k = 0$.

\section{Verification}\label{sec:verification}

The verification consists of three lemmas:
column-$0$ has gaps $\ge 2$ (Lemma~\ref{lem:gaps}), the construction
is a valid SYT (Lemma~\ref{lem:syt}), and there are no column
descents (Lemma~\ref{lem:no-col-descent}).

\begin{lemma}[Column-$0$ gaps]\label{lem:gaps}
The values $v_0 < v_1 < \cdots < v_{L+q-1}$ satisfy $v_0 = 1$,
$v_{L+q-1} = n$, and $v_{r+1} - v_r \ge 2$ for all $r$.
\end{lemma}

\begin{proof}
\textbf{If $k \ge 1$:} for $0 \le r < k-1$ we have $v_{r+1} - v_r = 2$.
At the boundary $r = k-1$,
\[
   v_k - v_{k-1} \;=\; \bigl(n - 2(L+q-1-k)\bigr) - (2k-1)
   \;=\; n - 2L - 2q + 3
   \;=\; M - q + 3 \;\ge\; 2,
\]
using $q \le M + 1$.  For $k \le r < L+q-1$ we have $v_{r+1} - v_r = 2$.
And $v_0 = 2 \cdot 0 + 1 = 1$, $v_{L+q-1} = n - 2 \cdot 0 = n$.

\textbf{If $k = 0$:} every row of $\hatl$ has length exactly $2$, so
$|\hatl| = 2L$, $M = 0$.  The constraint $\tau = 0$ becomes $q \le 1$,
hence $q = 1$ and $n = 2L+1$.  Then $v_r = n - 2(L+q-1-r) = 2r+1$,
giving $v_r = 1, 3, \ldots, 2L+1$.  All gaps are $2$.
\end{proof}

\begin{lemma}[Valid SYT]\label{lem:syt}
$T^*$ is a standard Young tableau of shape $\lambda$.
\end{lemma}

\begin{proof}
Every value in $\{1, \ldots, n\}$ is placed exactly once.  Row
strictness within column $0$: $v_r < v_{r+1}$ by
Lemma~\ref{lem:gaps}.  Column strictness within column $0$ is the
same statement.

\textbf{Strictness inside $\hatl'$.}  Within row $r$ of $\hatl'$, the
values $w_{a_r+1}, w_{a_r+2}, \ldots, w_{a_r + \hatl_r - 1}$ are
strictly increasing because the $w$ sequence is.  For column $c \ge
1$ of $\lambda$: the cell $(r, c)$ contains $w_{a_r + c}$, and the
cell $(r+1, c)$ (if in $\lambda$) contains $w_{a_{r+1} + c}$ with
$a_{r+1} = a_r + (\hatl_r - 1) \ge a_r + 1$.  Hence $w_{a_{r+1}+c} >
w_{a_r + c}$.

\textbf{Row condition across the boundary $c = 0 \to c = 1$.}  We
must show $T^*(r, 0) < T^*(r, 1)$ for all $r$ with $\hatl_r \ge 2$
(i.e.\ $0 \le r \le L-1$).  This is $v_r < w_{a_r + 1}$.

We treat the two ranges of $r$ separately.

\emph{Case $0 \le r < k$.}  Then $v_r = 2r+1$.  Also
\[
   a_r = \sum_{i < r}(\hatl_i - 1) \ge \sum_{i < r} 2 = 2r,
\]
since each long row $\hatl_i \ge 3$ for $i < r < k$.  Thus $a_r + 1
\ge 2r + 1$.  We split further.

If $a_r + 1 \le k - 1$, then $w_{a_r+1} = 2(a_r+1) \ge 2(2r+1) > 2r+1$.

If $a_r + 1 \ge k$, then $w_{a_r+1} \ge w_k = 2k$.  Since $r < k$,
we have $v_r = 2r + 1 \le 2k - 1 < 2k$.

\emph{Case $k \le r \le L-1$.}  Now $\hatl_r = 2$.  We compute $a_r$
explicitly:
\[
   a_r = \sum_{i < k}(\hatl_i - 1) + \sum_{k \le i < r}(\hatl_i - 1)
       = \bigl(|\hatl| - 2(L-k) - k\bigr) + (r - k)
       = M + r,
\]
using $\sum_{i<k} \hatl_i = |\hatl| - 2(L-k)$ and $\hatl_i = 2$ for
$i \ge k$.

We claim $w_{a_r + 1} = w_{M + r + 1}$ lies in Section S3.  Indeed
$M + r + 1 \ge M + k + 1 = (k + M - q + 1) + q > k + M - q + 1$,
i.e., past the end of Section S2.  Hence using \eqref{eq:sections}
with $i = M + r + 1$:
\[
   w_{M+r+1} = v_k + 2\bigl((M+r+1) - (k+M-q+1)\bigr) - 1
            = v_k + 2(r - k + q) - 1.
\]
Recall $v_k = n - 2(L+q-1-k)$, and $v_r = n - 2(L+q-1-r)$, so
\[
   v_k + 2(r - k + q) - 1 - v_r
   \;=\; \bigl(n - 2(L+q-1-k)\bigr) + 2(r-k+q) - 1 - \bigl(n - 2(L+q-1-r)\bigr)
\]
\[
   \;=\; -2(L+q-1-k) + 2(r-k+q) + 2(L+q-1-r) - 1
   \;=\; 2q - 1 \;\ge\; 1,
\]
using $q \ge 1$.  So $w_{a_r+1} - v_r \ge 1$ and $v_r < w_{a_r+1}$,
as required.

This completes the row condition, so $T^*$ is a valid SYT of shape
$\lambda$.
\end{proof}

\begin{lemma}[No column descents]\label{lem:no-col-descent}
$T^*$ has no column descent.
\end{lemma}

\begin{proof}
Suppose for contradiction that $j$ and $j+1$ occupy cells in the same
column at adjacent rows.  Let $(r, c)$ be the cell of $j$ and
$(r+1, c)$ the cell of $j+1$.

\textbf{Case A: $c = 0$.}  Both $j$ and $j+1$ are column-$0$ values,
so $j = v_r$ and $j+1 = v_{r+1}$.  But $v_{r+1} - v_r \ge 2$ by
Lemma~\ref{lem:gaps}, contradicting $j+1 - j = 1$.

\textbf{Case B: $c \ge 1$ and one of $j, j+1$ is in column~$0$.}
Impossible: $j$ and $j+1$ are in the same column $c \ge 1$, not column~$0$.

\textbf{Case C: $c \ge 1$, both $j, j+1$ in $\hatl'$.}  Then $j = w_i$
and $j+1 = w_{i'}$ for some $i < i'$.  Since $w$ is increasing and
$w_{i'} = w_i + 1$, we have $i' = i+1$ and $w_{i+1} - w_i = 1$.

In which Section is $i$?  Section S1 has gap $2$ between successive
terms.  Section S3 has gap $2$.  At the transition $i = k-1 \to k$
(end of S1 to start of S2) the values are $2k-2 \to 2k$ (gap $2$).
At $i = k+M-q+1 \to k+M-q+2$ (end of S2 to start of S3) the values
are $(2k + M - q + 1) \to (v_k + 1) = (n - 2(L+q-1-k) + 1)$;
using $n = 2L + M + q$, $v_k + 1 = 2k + M - q + 3$, so the gap is
$2$.  The only place with gap $1$ is \emph{inside} Section S2, where
$w_{i+1} - w_i = 1$.

So $i$ and $i+1$ both lie in S2, i.e.,
\begin{equation}\label{eq:S2-range}
   k \;\le\; i \;<\; i+1 \;\le\; k + M - q + 1.
\end{equation}

Now $w_i$ and $w_{i+1}$ lie in adjacent rows $r, r+1$ of $\hatl'$ at
the same column~$c$.  In row-superstandard order, $w_i$ is the last
cell of row~$r$ iff $i = a_{r+1}$ (since row $r$ contains the indices
$a_r + 1, \ldots, a_r + (\hatl_r - 1) = a_{r+1}$), and then $w_{i+1}$
is the first cell of row $r+1$.  The column of $w_i$ within $\hatl'$
is $\hatl_r - 1 - 1 = \hatl_r - 2$ (the last column of row $r$ of
$\hatl'$), while the column of $w_{i+1}$ within $\hatl'$ is $0$.
Translating to columns of $\lambda$ (add $1$): $c = \hatl_r - 1$ for
$w_i$, and $c = 1$ for $w_{i+1}$.  These match iff $\hatl_r = 2$ and
the column is $c = 1$.

Thus $r \ge k$ (so that $\hatl_r = 2$) and $r \le L-1$.  If $r = L-1$,
the cell directly below $(L-1, 1)$ is $(L, 1)$, but $\lambda_L \le 1$
(we are in the tail), so $(L, 1) \notin \lambda$, contradicting our
hypothesis that $(r+1, c) = (L, 1)$ is filled.  Hence $r \le L-2$.

For $r \in [k, L-2]$ we have $\hatl_r = \hatl_{r+1} = 2$ and (from
the proof of Lemma~\ref{lem:syt}) $a_{r} = M + r$ for $r \ge k$, so
$a_{r+1} = M + r + 1$.  Hence $i = a_{r+1} = M + r + 1$, and
\eqref{eq:S2-range} becomes
\[
   k \;\le\; M + r + 1 \quad\text{and}\quad M + r + 2 \;\le\; k + M - q + 1,
\]
i.e., $r \ge k - M - 1$ (trivially satisfied since $r \ge k$ and
$M \ge 0$) and $r \le k - q - 1$.  But $r \ge k$ and $q \ge 1$ forces
$k - q - 1 < k \le r$, a contradiction.

Therefore no column descent exists.
\end{proof}

\begin{proof}[Proof of Theorem~\ref{thm:main}]
By Lemma~\ref{lem:syt}, $T^*$ is a valid SYT of shape $\lambda$.  By
construction $T^*(0,0) = v_0 = 1$.  For $T^*(0,1) = w_1$: if $k \ge 2$,
$w_1 = 2$ from S1; if $k = 1$, $w_1 = 1 + k = 2$ from S2; if $k = 0$,
$w_1 = 2 \cdot 1 = 2$ from S3 (using $v_0 = 1$).  In all cases
$T^*(0,1) = 2$.  By Lemma~\ref{lem:no-col-descent}, $T^*$ is
column-descent-free.
\end{proof}

\section{Computational verification}\label{sec:computation}

The construction $T^*$ has been implemented and verified to produce
column-descent-free SYTs on every $\tau = 0$ shape with $|\lambda|
\le 25$:

\begin{quote}
\verb|/home/clio/projects/scratch/2026-05-16-prove-non-stuck/t_unified.py|.
\end{quote}

\begin{itemize}
\item $|\lambda| \le 14$: $210/210$ shapes pass.
\item $|\lambda| \le 18$: $832/832$ shapes pass.
\item $|\lambda| \le 22$: $2688/2688$ shapes pass.
\item $|\lambda| \le 25$: $5957/5957$ shapes pass.
\end{itemize}

\section{Remarks}\label{sec:remarks}

\subsection{Comparison with the greedy Algorithm~N}

The companion paper \cite{TstarGreedy} introduced a greedy algorithm
(Algorithm~N) whose non-stuckness conjecture (Conjecture~4.18)
reduces to Conjecture~4.6.  The present construction $T^*$ is
\emph{not} the output of Algorithm~N in general; it is a different,
fully explicit deterministic recipe.  Either suffices to close the
existence question.

\subsection{Where the structure ``hides''}

The proof identifies a single arithmetic fact as the source of
column-descent-freeness: in Section S2 of the value sequence
$w_1 < w_2 < \cdots$, the only place where consecutive integers
occur, the index range $[k, k+M-q+1]$ ends \emph{strictly before}
the index $M + r + 1$ corresponding to the first cell of row
$r+1 \ge k+1$ of $\hatl'$.  The gap is exactly $q \ge 1$ indices.
Lose $q$ and the construction breaks; the $\tau = 0$ hypothesis
is exactly the amount of slack needed.

\subsection{Closing the trace-vanishing equivalence}

Combining Theorem~\ref{thm:main} with the converse direction of the
per-SYT positivity proof in \cite{ConversePerSYT} (which assumed the
existence of $T^*$), and with the forward direction of trace-vanishing
already proved in \cite{TraceForward}, we conclude:

\begin{corollary}[Trace-vanishing $\Leftrightarrow$ $\tau \ge 1$]
For $\lambda = (\hatl, 1^q)$ with $\hatl_r \ge 2$ for all $r$ and
$L \ge 2$, the trace $\mathrm{tr}_{V^\lambda}(\Omega^{(\lambda)})
\in \mathbb{Q}_{\ge 0}(q)$ vanishes identically as a rational function
of $q$ if and only if $\tau(\hatl) \ge 1$.
\end{corollary}

\begin{thebibliography}{9}
\bibitem{ConversePerSYT}
Clio, \emph{Per-SYT positivity of $\Omega^{(\lambda)}$ in the asymmetric
Hoefsmit basis}, May~2026 (commit \texttt{3282849}).

\bibitem{TstarGreedy}
Clio, \emph{A deterministic column-balanced greedy construction for the
column-descent-free SYT existence conjecture}, May~2026 (commit
\texttt{6f17ed8}).

\bibitem{TraceForward}
Clio, \emph{Trace-vanishing of $\Omega^{(\lambda)}$ from the multi-row
sign-kill meta-theorem}, May~2026 (commit \texttt{92a5323}).
\end{thebibliography}

\end{document}
